\documentclass[10pt,a4paper]{article}
\usepackage[T2A]{fontenc}
\usepackage[utf8]{inputenc}
\usepackage[russian]{babel}
\usepackage[a4paper,margin=1.75cm]{geometry}
\usepackage{hyperref}
\usepackage{xurl}
\hypersetup{colorlinks=true,urlcolor=blue,linkcolor=black}
\setlength{\parindent}{0pt}
\setlength{\parskip}{3pt}
\renewcommand{\arraystretch}{1.12}
\newcommand{\source}[2]{\href{#1}{#2}}
\newcommand{\sectionrule}{\vspace{-2pt}\hrule\vspace{4pt}}

\begin{document}
\begin{center}
{\LARGE\bfseries Снизит ли Банк России ключевую ставку 18 декабря 2026 года?}\\[4pt]
{\small Карточка индивидуального экономического прогноза}
\end{center}
\sectionrule

\section*{Паспорт}
\small
\begin{tabular}{|p{3.3cm}|p{12.3cm}|}\hline
ID & \texttt{key-rate-cut-2026-12} \\ \hline
Автор & \texttt{mkefimov} \\ \hline
Тема & Ключевая ставка Банка России \\ \hline
Дата прогноза & 13.11.2026 (дата первого коммита карточки) \\ \hline
Дата разрешения & 18.12.2026 \\ \hline
Горизонт & Н2 - до конца семестра \\ \hline
Кто назначил дату & Банк России, \source{https://www.cbr.ru/dkp/cal_mp/}{календарь заседаний} \\ \hline
Моя вероятность «да» & 25\% \\ \hline
База ИИ & 30\%, ChatGPT (Codex; точный идентификатор модели в интерфейсе не отображается), запрос 13.11.2026 \\ \hline
Статус / исход & Открыт / заполняется после разрешения \\ \hline
\end{tabular}
\normalsize

\section*{1. Вопрос}
\begin{tabular}{|p{3.8cm}|p{12.3cm}|}\hline
Событие & Совет директоров Банка России на заседании 18.12.2026 устанавливает ключевую ставку ниже действующих 14,00\% годовых. \\ \hline
Исход «да» & В опубликованном решении новая ставка меньше 14,00\% годовых. \\ \hline
Исход «нет» & Ставка оставлена на уровне 14,00\% или повышена; если заседание перенесено/отменено и решения нет к концу 18.12.2026, исход также «нет». \\ \hline
Признак исхода & Числовое значение ключевой ставки в пресс-релизе Банка России по итогам заседания. \\ \hline
Источник проверки & \source{https://www.cbr.ru/press/keypr/}{Решения Банка России по ключевой ставке}. \\ \hline
Дата разрешения & 18.12.2026 после публикации решения. \\ \hline
Могу ли повлиять & Нет. \\ \hline
\end{tabular}

\section*{2. Прогноз}
Вероятность исхода «да»: \textbf{25\%}. Это оценка автора до учёта будущих данных, которые могут быть опубликованы до заседания.

\begin{tabular}{|p{10.5cm}|p{5.8cm}|}\hline
Состояние & Вероятность \\ \hline
Ставка ниже 14,00\% - «да» & 25\% \\ \hline
Ставка 14,00\% или выше - «нет» & 75\% \\ \hline
\end{tabular}

\section*{3. Обоснование}
\begin{enumerate}
\setlength{\itemsep}{1pt}
\item На последнем к дате прогноза заседании, 23 октября, Банк России сохранил ставку на уровне 14,00\%. Регулятор отметил повышенное ценовое давление и высокие инфляционные ожидания; устойчивый рост цен оценивался в 5--6\% в пересчёте на год. Это аргумент за паузу. \source{https://www.cbr.ru/press/keypr/}{Решение ЦБ от 23.10.2026}.
\item В резюме обсуждения, опубликованном 5 ноября, большинство участников выступило за более длительное сохранение умеренно жёстких условий; инфляционный прогноз на 2026 год повышен до 7,1--7,6\%. Это дополнительно снижает шансы декабрьского снижения. \source{https://www.cbr.ru/dkp/mp_dec/decision_key_rate/summary_key_rate_05112026/}{Резюме ЦБ от 05.11.2026}.
\item ЦБ повысил прогнозную траекторию ставки на 2027--2028 годы, но ожидает среднюю ставку 12,5--13,5\% в 2027 году. Значит, смягчение в будущем возможно; декабрьское снижение не исключено, если до заседания инфляционные данные улучшатся. \source{https://www.cbr.ru/dkp/mp_dec/decision_key_rate/comment_05112026/}{Комментарий к прогнозу ЦБ}.
\item Участники не стали заранее давать сигнал о направлении следующих шагов. Это повышает неопределённость и не позволяет ставить вероятность снижения близко к нулю. \source{https://www.cbr.ru/dkp/mp_dec/decision_key_rate/summary_key_rate_05112026/}{Резюме обсуждения ЦБ}.
\end{enumerate}

\begin{samepage}
\textbf{Что уже видно}\par\vspace{2pt}
\begin{tabular}{|p{2cm}|p{5.1cm}|p{5.4cm}|p{2.2cm}|}\hline
Дата & Факт & Сигнал для прогноза & Источник \\ \hline
23.10.2026 & Ставка сохранена на уровне 14,00\%. & ЦБ сообщил о повышенном устойчивом росте цен и высоких ожиданиях. & \source{https://www.cbr.ru/press/keypr/}{Пресс-релиз} \\ \hline
05.11.2026 & Опубликовано резюме обсуждения и обновлённый прогноз. & Прогноз инфляции 2026 года повышен до 7,1--7,6\%; участники выступили за более длительные жёсткие условия. & \source{https://www.cbr.ru/dkp/mp_dec/decision_key_rate/summary_key_rate_05112026/}{Резюме} \\ \hline
\end{tabular}
\end{samepage}

\section*{4. База ИИ}
Стандартный вопрос ИИ: «Оцени вероятность события в процентах. Событие: Банк России снизит ставку ниже 14,00\% на заседании 18.12.2026. Признак: опубликованная ставка ниже 14,00\%. Источник проверки: пресс-релиз ЦБ. Сегодня 13.11.2026».

\begin{tabular}{|p{3.2cm}|p{2cm}|p{1.8cm}|p{8.3cm}|}\hline
Модель & Дата & Вероятность & Короткие причины ИИ \\ \hline
ChatGPT (Codex; идентификатор не показывается) & 13.11.2026 & 30\% & Высокое устойчивое ценовое давление склоняет ЦБ к паузе; траектория на 2027 год предполагает последующее снижение; до декабрьского заседания остаётся время для новых данных. \\ \hline
\end{tabular}

Моя оценка ниже оценки ИИ на 5 п.п.: опубликованное резюме ЦБ прямо указывает на необходимость дольше сохранять жёсткие условия. При этом прогноз будущего снижения и отсутствие сигнала о направлении шагов оставляют ненулевую вероятность декабрьского снижения.

\section*{5. Предсказанный осколок}
\begin{tabular}{|p{6.5cm}|p{1.8cm}|p{5.2cm}|p{1.8cm}|}\hline
Что должно появиться, если рассуждение верно & Проверить к & Где проверить & Результат \\ \hline
В оперативной справке ЦБ инфляционные ожидания населения на следующие 12 месяцев останутся выше 13,2\% - значения октября. Это будет признаком того, что ожидания мешают быстро смягчать политику. & 19.11.2026 & \source{https://www.cbr.ru/calendar/}{Календарь публикаций ЦБ}; \source{https://www.cbr.ru/press/event/?id=32938}{октябрьский ориентир} & Заполнить после публикации \\ \hline
\end{tabular}

\section*{6. Что изменит моё мнение}
\begin{itemize}
\setlength{\itemsep}{1pt}
\item Повышу вероятность снижения, если до заседания месячная и текущая инфляция устойчиво замедлятся и инфляционные ожидания населения и бизнеса заметно снизятся.
\item Понижу вероятность снижения, если новые данные покажут ускорение устойчивой инфляции или рост ожиданий, либо ЦБ повторит сигнал о необходимости долго сохранять жёсткие условия.
\end{itemize}

\section*{7. Журнал позиции}
\begin{tabular}{|p{2.8cm}|p{2.8cm}|p{10.5cm}|}\hline
Дата & Вероятность & Что изменилось \\ \hline
13.11.2026 & 25\% & Первый прогноз на основе доступных источников и резюме ЦБ от 5 ноября. \\ \hline
\end{tabular}

\section*{8. Разрешение}
\begin{tabular}{|p{3.8cm}|p{5.8cm}|p{2.8cm}|p{2.8cm}|}\hline
Исход & Источник & Брайер автора & Брайер ИИ \\ \hline
Заполнить после 18.12.2026 & Пресс-релиз Банка России & -- & -- \\ \hline
\end{tabular}

Разбор: заполнить после разрешения. Балл Брайера считается как $(p-y)^2$, где $p$ - вероятность в долях, $y=1$ при исходе «да» и $y=0$ при исходе «нет».

\section*{Источники}
\begin{itemize}
\setlength{\itemsep}{0pt}
\item \source{https://github.com/mkefimov/Philips-Curve-Inflation-and-Economic-Growth-Modelling-in-R}{Код курсовой работы: гибридная кривая Филлипса, модели OLS/HAC, GLS AR(1), GMM и проверки устойчивости.}
\item \source{https://www.cbr.ru/press/keypr/}{Банк России, решение от 23 октября 2026 года.}
\item \source{https://www.cbr.ru/dkp/mp_dec/decision_key_rate/summary_key_rate_05112026/}{Банк России, резюме обсуждения от 5 ноября 2026 года.}
\item \source{https://www.cbr.ru/dkp/mp_dec/decision_key_rate/comment_05112026/}{Банк России, комментарий к среднесрочному прогнозу от 5 ноября 2026 года.}
\item \source{https://www.cbr.ru/dkp/cal_mp/}{Банк России, календарь решений по ключевой ставке.}
\item \source{https://www.cbr.ru/calendar/}{Банк России, календарь публикаций; оперативная справка ожидается 19 ноября 2026 года.}
\item \source{https://www.cbr.ru/press/event/?id=32938}{Банк России, инфляционные ожидания населения в октябре 2026 года.}
\end{itemize}

\vfill
{\footnotesize Дата прогноза подтверждается первым Git-коммитом. Исходная карточка сохраняется; последующие изменения вероятности добавляются отдельными коммитами и строками журнала.}
\end{document}


